\section{Conclusion}

Two things can be wrong with a trustworthy forecasting system without any of its reported numbers looking wrong. It can ignore most of the input it is configured to read, and the score it uses to judge its own forecasts can be worse than one of the signals it is built from. This paper contributes a diagnostic for each, and a design corrected by them.

The effective-receptive-field audit is a training-free perturbation sweep that reports how much of its history a model can actually condition on, measured on the reported weights and gated on their reproducing the ledger. It found a published temporal stem admitting \robustFieldPublished{} of \robustHistory{} configured steps, a gap that no accuracy, coverage, or interpretability measurement in an earlier full campaign had surfaced. Replacing that stem with a causal dilated stack whose reach is a closed-form property of its depth, and adding an attention readout so a year of context need not pass through one final state, gives a model that reads its window and, in the paired family, is significantly better than a budget-matched plain KAN, better than the architecture it replaces at every horizon, and better than a wide-stem-only variant at \pairedDilatedWinHorizons{} and nowhere else.

Against the strongest baseline the corrected model is statistically tied: \pairedTransformerWins{} of \pairedTransformerTotal{} pairs favour it, \pairedTransformerLosses{} favours the transformer, and \pairedTransformerTied{} are inconclusive after Holm correction. We claim no accuracy advantage there and have written the generators so the tables cannot imply one.

Calibration-selected fusion replaces a frozen combination weight with a frozen selection procedure, over a grid whose endpoints are the single-component scores. The construction removes the fusion-loses-to-its-own-component failure mode outright, and it allows the data to withhold mass from an unhelpful signal: on CET the procedure put \fusionWeightLow{}--\fusionWeightHigh{} of the mass on interval width and never selected equal weights in \fusionRuns{} runs. The reliability behaviour that follows is the substance of the paper. Error detection rises from \errorAurocOne{} at one day to \errorAurocNinety{} at ninety, simultaneous joint coverage holds at \simulCoverageThirty{}--\simulCoverageNinety{} against a nominal \nominalCoverage{}, and rank correlation with realized error reaches $|\rho|=\reliabilityRhoMax{}$.

The boundaries are as much a result as the gains. Day-ahead error detection is at chance because day-ahead error on this record is irreducible from anything internal to the model. The shift term's marginal contribution to selective risk stays negligible even once its weight is fitted; what the fitting changed is that the method reaches that conclusion itself rather than requiring a refutation to establish it. The correction costs several times the training time of the cheapest neural baseline. And all of this is one station: the station-generalization protocol is frozen and unexecuted, and it is the obvious next obligation.

The general lesson concerns what reliability evidence has to be about. A configured context length is not an effective one, and only the second is a property of the model; a frozen constant is not a conservative choice when it governs how evidence is combined, and only a frozen procedure is. Both distinctions cost almost nothing to check, and in this study each of them was the difference between a system that reported reliability and one that had it.
